\hypertarget{5-discussion}{%
\section{Discussion}\label{5-discussion}}

\hypertarget{51-interpreting-track-a}{%
\subsection{Interpreting the Track A negative result}\label{51-interpreting-track-a}}

The Track A result should be read as evidence of a structural limitation
rather than evidence that RL is ineffective for supply chains in general.
Under the thesis-grounded contract, the static S2 policy already operates
near the practical ceiling imposed by downstream distribution capacity:
Op10--Op12 transport throughput of approximately 2,500 rations per day
roughly matches demand, so additional upstream production capacity from S3
generates excess inventory in intermediate buffers without proportionally
increasing deliveries. A design-of-experiments analysis confirms an
asymmetric action landscape: destructive inventory settings reduce reward by
44--65\%, while positive headroom above S2-neutral is limited to
approximately 1\%. PPO's convergence to near-neutral inventory actions with
S2-region shift selection is therefore rational risk management rather than
learning failure. Reward redesign (14 modes tested), memory (RecurrentPPO),
and richer observations (v2--v6) do not overturn this structural limitation
because the binding constraint lies outside the controllable interface.

\hypertarget{52-what-track-b-changes}{%
\subsection{What Track B changes}\label{52-what-track-b-changes}}

Track B changes the controllable interface, not the underlying DES physics.
By adding Op10 and Op12 dispatch multipliers to the action space, the agent
gains leverage over the downstream bottleneck that constrains Track A. The
learned policy achieves higher Excel ReT, service continuity, and
recovery-tail metrics than the best of a dense downstream-dispatch grid of static comparators,
but it is not a strict cost win on the shift dimension alone: its
shift-utilization cost index is not statistically distinguishable from
the best comparator selected by ReT at zero dispatch charge. We report
the result as Pareto non-dominated on the ReT--cost--tail--flow frontier. The
dispatch-inclusive cost sensitivity (Section~4.3) sharpens the managerial
reading: because the learned policy expedites selectively rather than
holding aggressive dispatch permanently, PPO becomes cheaper in total than
the best static once dispatch expediting carries a positive charge at the
tested threshold --- the policy buys recovery with \emph{targeted}, not free,
transport flexibility.

The action-space decomposition (Section~4.4) strengthens this
interpretation. At
matched 5-seed, 60{,}000-timestep scale, the \emph{shift\_only} condition
achieves a comparable advantage over its own best evaluated comparator
(+0.000377) to the full \emph{joint} contract (+0.000367), while
\emph{downstream\_only} exceeds both (+0.000429). This decomposition
provides mechanism evidence consistent with downstream dispatch access
being the strongest observed lever and with the number of controllable
dimensions not being what drives the Track B
advantage; it does not establish downstream dispatch as the only useful
control surface, since shift-only control also captures part of the
dynamic effect.

\hypertarget{53-action-space-alignment}{%
\subsection{Action-space alignment as a design principle}\label{53-action-space-alignment}}

The dual-track result connects three literatures that have not previously
been linked empirically: the Theory of Constraints \cite{goldratt1984goal},
which identifies the binding constraint as the lever that determines system
throughput; organizational-learning accounts of adaptation and forgetting
\cite{levitt1988organizational}; RL agent design, which determines what the
agent can control; and action-space architecture, which determines how
decisions are represented.
Our benchmark demonstrates a concrete instance of this connection: when the
action space does not cover the active constraint (Track A), better
algorithms and rewards do not convert the measured headroom; when downstream
dispatch authority is exposed (Track B), a standard PPO+MLP agent can produce
a decisive advantage in the designed sustained-disruption cell. The reward
sweep (Section~4.6) reinforces this: all 18 screened reward/observation cells
show positive Excel ReT deltas, whereas in Track A no reward formulation
succeeds. This suggests a practical design principle: \emph{before investing
in architectural novelty or reward engineering, verify that the agent's
controllable interface reaches the binding operational bottleneck and that
the resulting headroom is dynamically convertible into recovery performance.}

\hypertarget{54-positioning}{%
\subsection{Positioning relative to the AI--SCRES literature}\label{54-positioning}}

Recent work at the intersection of AI and supply chain resilience spans
several paradigms. Rajagopal and Sivasakthivel~\cite{rajagopal2024} use weighted
feed-forward neural networks for resilience \emph{strategy selection}---a
classification task fundamentally different from our online adaptive control.
Rezki and Mansouri~\cite{rezki2024} apply neural networks to \emph{predict} supplier
risk, while Ordibazar et al.~\cite{ordibazar2022} use counterfactual explanations for
risk \emph{recommendation}. These contributions address complementary
problems (prediction, classification, explanation) rather than the
prescriptive operational control explored here. Kogler and Maxera~\cite{kogler2026}
review 72 papers combining DES with ML for supply chain analysis, in
which DES+RL emerges as the most common hybrid pairing; our work contributes
a specific instance with validated thesis-grounded data, honest negative
results, and mechanism-ablation evidence that is rare in this literature.

Among RL-specific competitors, Ding et al.~\cite{ding2026} apply MAPPO for dynamic
reconfiguration of supply chain topology under disruption, operating at the
strategic level with abstract network actions. Our single-agent approach
asks a complementary question at the operational level: given a fixed
topology, does expanding operational control scope change the outcome? The
Track B result suggests that, in a fixed-topology supply chain where the
binding downstream bottleneck can be identified, centralized control with
appropriate action scope can be a strong benchmark before moving to more
complex multi-agent architectures.

\hypertarget{55-limitations}{%
\subsection{Limitations}\label{55-limitations}}

Several limitations bound the generalizability of these findings. First, the
benchmark is one validated military supply chain; generalization to other
topologies, demand patterns, or organizational structures requires further
investigation. Second, Track B is not a thesis-grounded decision-variable
replication of Garrido-Rios: it preserves the DES backbone and resilience
measurement structure, but extends the controllable surface to downstream
dispatch valves that the original experiment held fixed. Third, the
observation ablation removes the seven explicitly privileged regime/forecast
fields, but the remaining state still contains rich operational feedback; the
experiment rules out the strongest privileged observation confound but does
not prove full field deployability of every state variable.
Fourth, the problem is better described as partially observed control than a
fully identified MDP: the agent observes a projection of the simulation state,
not the full internal state. Fifth, the paper uses PPO exclusively (with
RecurrentPPO as a comparator) as a stable continuous-control baseline; we
claim action-space alignment, not algorithmic superiority, and off-policy
continuous-control algorithms such as SAC and TD3 are natural robustness
extensions. Algorithmic choice is not the mechanism tested in this paper,
and the reward-insensitivity screen suggests the effect is not tied to
one training configuration. Sixth, the
cross-regime/horizon matrix uses a fixed dense-comparable static comparator
set rather than a fully reoptimized 147-cell dense frontier in every
risk/horizon cell. Seventh, the added 3$\times$3 upstream bound tests
nearby Op3/Op9 static quantity choices at the best downstream cell and
does not overturn PPO, but it is still a local bound rather than an
exhaustive 8D static frontier. Finally, the evidence supports adaptive recovery and
backlog control, not unqualified anticipation, cross-campaign organizational
learning, universal cost reduction, or a universal law that bottleneck access
alone guarantees RL success. A separate, confirmatory-scale probe
(10 seeds, 8 sequential online-adaptation cycles) finds a small positive
retained-versus-reset advantage in Excel ReT ($+0.0000493$, 95\% CI
$[+0.0000167, +0.0000819]$, masked-observation arm), roughly an order of
magnitude smaller than the Track B headline gain and robust to masking
the same privileged regime/forecast fields removed in
Section~\ref{45-privileged-obs}; this is preliminary evidence for a
retained-adaptation channel distinct from the action-space alignment
mechanism studied here, not a claim of cross-campaign organizational
learning, and is not part of the paper's central result.

\hypertarget{56-future-work}{%
\subsection{Future work}\label{56-future-work}}

Three directions emerge. First, the bottleneck-aligned recovery principle
should be tested on other validated supply chain DES models to assess
generalizability beyond the MFSC topology. Second, more expressive policy
architectures---including attention-based networks, graph neural networks,
and Kolmogorov--Arnold networks, which Garrido et
al.~\cite{garrido2024iccl} name (alongside backpropagation networks and
reinforcement learning) as candidate families for closing the
learning gap in SCRES-DES models, though without ranking them or
arguing for KAN-specific properties---remain relevant for supply chains with variable topology or
multi-agent coordination, but were not required for this paper's result: a
standard PPO+MLP architecture already produces a robust, adversarially-tested
advantage under Track B. Third, the generalization matrix (Section~4.7)
identifies severe disruption at both evaluated horizons as a genuine boundary case;
training specialized policies for that regime, rather than relying on the
policy trained under \texttt{adaptive\_benchmark\_v2}, could establish
tighter upper bounds on RL effectiveness under catastrophic disruption.
Cross-campaign organizational learning at longer horizons than the
8-cycle probe reported above, off-policy algorithm comparisons
(SAC/TD3), and a full 147-candidate regime-table fit (the current
privileged-observation test used a trimmed 75-candidate grid) remain open
for future work.
